\scanpage{214}
\mspage{451}
% archive sheet number: 213
résiste pas à la tentation de la reproduire ici \inserted{(traduite)} in extenso.\deleted{\handwritten{\unsure{-}}} Il m'a fallu un moment, d'ailleurs, pour saisir que c'\deleted{\unsure{e}}était bel et bien une réponse à ma belle lettre \deleted{\unsure{d'il y avait}} du mois dernier\deleted{\unsure{,}}. \inserted{\handwritten{Voici donc la réponse.}}

Heidelberg 15.3.1985\\
\deleted{\handwritten{\unsure{\char34{}}}}Cher Professeur Grothendieck,\\

Il me faut encore vous remercier pour votre lettre du 9 février. La lettre de Mme Dr. Byrne du 15 février aura sans doute répondu à vos questions.

Rece\deleted{g}vez e\deleted{\unsure{s}}tc\\
K.F. Springer \deleted{\handwritten{\unsure{\char34{}}}}\\

Là au moins je suis fixé ! Les gens \char34{}bien informés\char34{} (dont il a été déjà question) ont dû lui expliquer que c'était pas la peine qu'il se fatigue pour le \deleted{\unsure{m}}Monsieur un peu excité qui lui écrivait là - qu'i\deleted{\unsure{b}}l ne faisait décidément pas partie du beau monde. Et c'est vrai, en plus ...

En attendant de recevoir cette réponse édifiante de la direction de l'entreprise de \deleted{\unsure{p}}Pompes \deleted{\unsure{f}}Funèbres Springer Verlag GmbH (c'était gentil encore de m'honorer d'une réponse \deleted{personelle}\inserted{signée} du directeur en personne), j'ai eu le temps de me sonder sur mes propres intentions. Le rôle joué par l'estimable entreprise me paraît vraiment très gros, et j'ai songé à l'éventualité d'un procès à gránd spectacle, où je demanderais des dommages-intérêts astronomiques, à titre de \char34{}monsieur bien\char34{}\handwritten{/}outragé, \deleted{et} victime d'inqualifiables passe\deleted{s}-droits . Mais je me suis dit aussi \deleted{que de tels pro} qu'un procès c\deleted{\unsure{d}}omme ça, ça doit bouffer une énergie\handwritten{\(\vert\)}dingue. A supposer même que j'aie gain \deleted{\unsure{d}}de cau\deleted{\unsure{d}}se et que je touche des dommages-intérêts vertigineux (soyons optimistes\unsure{)|!})\handwritten{,} au bout de X années\deleted{\unsure{,}}\inserted{\handwritten{\unsure{certes|cuites}}} - à quoi serais-je avancé ? Je ne suis pas dans le besoin et n'ai nul besoin de plus que ce que j'ai - et une escroquerie n'est pas plus ni moins une escroquerie, parce qu'un certain procès a été gagné, ou perdu. Je ne vais pas amél\deleted{\unsure{oi}}iorer le monde, ni moi-même, ni les manières de Monsieur K.F. Springer et de certains employés de l'entreprise qu'il dirige, et en tous cas pas leur\deleted{s} façon de concevoir leur métier\handwritten{,} en mobilisant des avocats et en \deleted{\unsure{leu}}\handwritten{les} faisant mobiliser les leurs\inserted{\handwritten{(*)}}. \deleted{Par contre, dans les années qui m} Et je n'améliorerai pas non plus un certain esprit dans un certain beau monde que j'ai quitté, \deleted{un}\handwritten{l'}esprit qui rend possible le genre d'opération dont le Dr. Springer et son estimable \deleted{entre} maison se \deleted{f}\handwritten{s}ont \inserted{\handwritten{faits depuis treize ans \deleted{\unsure{déjà}}}} \deleted{\unsure{l}}\handwritten{l}es serviteurs. Il me reste (je l'espère bien) quelques
